\begin{textsample}{POS Dim 1 – vem\_ed\_24 – Score 15.00 – vem\_ed\_24/t129.txt}  \label{ex:f1_pos_055}
IVES Cesu 2024: Inteligência Artificial e inovação Gabriela Méndez, da Florida International University (EUA) Gabriela Méndez (Florida International University, EUA) exemplificou \textbf{estratégias} \textbf{práticas} para inovação em projetos COIL \textbf{utilizando} inteligência \textbf{artificial}.
Sua apresentação está disponível em http://go.fiu.edu/ives2024 A Inteligência Artificial Generativa funciona mediante o uso de um modelo de aprendizagem de máquina (machine learning).
Essa aprendizagem ocorre a \textbf{partir} de \textbf{conteúdos} \textbf{criados} por \textbf{humanos} (conjuntos de dados, padrões, relações).
A IA Generativa \textbf{cria} novos \textbf{conteúdos} a \textbf{partir} dos padrões aprendidos, \textbf{é} o caso dos chatbots (ChatGPT, Microsoft Copilot, Perplexity, Koala Chat e You.com), dos geradores de imagens (playground.com, Leonardo.ai, Stable Diffusion On-line), de apresentações (gamma.app), de vídeos (fliki.ai, genmo.ai) e de voz (elevenlabs.io).
Para interagir com a IA Generativa, \textbf{é} preciso redigir o prompt (pergunta, pedido, instrução) de maneira clara e \textbf{específica}, para obter respostas mais adequadas ao \textbf{que} estamos buscando. \textbf{É} preciso incluir o \textbf{contexto}, adicionar detalhes, exemplos ou cenários \textbf{relevantes} para garantir \textbf{que} a mensagem se \textbf{alinhe} com seus objetivos. “Experimente e aperfeiçoe, pratique até \textbf{que} obtenha resultados com os quais esteja satisfeito”, recomenda Gabriela.
A IA Generativa \textbf{é} útil em projetos COIL, especialmente para brainstorming, design instrucional, comunicação e matchmaking entre professores de áreas \textbf{que} à primeira vista podem parecer distantes.
O Chat GPT pode ajudar a desenvolver ideias para projetos COIL, por exemplo. “Daqui a cinco anos a educação vai \textbf{ser} mais autônoma, mas os estudantes precisarão de alguém para guiar seus passos na instituição de ensino”, projeta Gabi, como \textbf{é} conhecida no mundo COIL. “Os alunos \textbf{serão} mais autodidatas, mas sempre com \textbf{apoio} do docente, \textbf{que} deve ter um \textbf{papel} de coach, de facilitador, de guia; para isso, temos \textbf{que} transformar nosso \textbf{papel} de educadores”.
Sobre o espírito questionador, imprescindível nas relações entre as inteligências \textbf{humana} e \textbf{artificial}, levantado na sessão de perguntas após a apresentação, Gabi respondeu: “\textbf{Todos} temos responsabilidade de tratar criticamente as ferramentas de IA”.
Se as \textbf{referências} da IA não \textbf{representam} os diferentes \textbf{contextos} sociais, \textbf{é} importante trazer visibilidade para esses \textbf{contextos} e, “para \textbf{que} isso aconteça, \textbf{é} preciso colaborar, \textbf{utilizando} essas ferramentas”.
Osvaldo Succi Junior finalizou as discussões sobre o tema com uma importante observação: “Especialmente nós \textbf{que} estamos no Sul Global temos \textbf{que} pensar muito cuidadosamente sobre essa \textbf{questão} da representatividade”.
Usar as ferramentas de IA com criticismo, sempre.

% matched lemmas: alinhar, apoio, artificial, contexto, conteúdo, criar, específico, estratégia, humano, papel, partir, prático, que, questão, referência, relevante, representar, ser, todo, utilizar
\end{textsample}
